\section{Activation age and lazy Schur settlement}
\label{sec:lazy-schur}

The revisit factor admits a dual expression that makes the fan obstruction
transparent.  For a nested trace
$U_0\subseteq\cdots\subseteq U_K$, define the activation time
\begin{equation}
 \tau(v):=\min\{k:v\in U_k\}
 \qquad(v\in U_K).
 \label{eq:activation-time}
\end{equation}

\begin{proposition}[Activation-age identity]
\label{prop:activation-age}
Every nested trace satisfies
\begin{align}
 \sum_{k=0}^K\operatorname{cvol}(U_k)
 &=
 \sum_{v\in U_K}(1+d_v)\bigl(K-\tau(v)+1\bigr),
 \label{eq:activation-age-sum}\\
 \mathfrak R_{\rm set}
 &=
 \sum_{v\in U_K}
 \frac{1+d_v}{\operatorname{cvol}(U_K)}
 \bigl(K-\tau(v)+1\bigr).
 \label{eq:activation-age-mean}
\end{align}
Thus $\mathfrak R_{\rm set}$ is the final-charge-weighted mean number of
fresh settlements paid after activation.
\end{proposition}

\begin{proof}
A fixed $v\in U_K$ appears in exactly the regions
$U_{\tau(v)},\ldots,U_K$.  Exchange the order of summation in
$\sum_k\sum_{v\in U_k}(1+d_v)$ and divide by
$\operatorname{cvol}(U_K)$.
\end{proof}

For the fan, the root has activation time zero and charge $m+1$, a constant
fraction of the full graph charge.  Its long lifetime alone yields
\cref{eq:fan-revisit-lower}.  An incremental method should therefore avoid
charging persistent coordinates merely because a new coordinate entered.
The exact algebraic interface is the following block update.

\begin{lemma}[Lazy block-Schur settlement]
\label{lem:lazy-block-schur}
Let $x$ be the shifted Dirichlet solution on $U$, with $x_{U^c}=0$, and let
$B\subseteq U^c$ be a nonempty block.  Put
\begin{equation}
 A:=Q_{UU},
 \qquad E:=Q_{UB},
 \qquad
 \mathcal S_B:=Q_{BB}-E^TA^{-1}E,
 \qquad
 q_B:=e_\ell(x)_B.
 \label{eq:lazy-schur-data}
\end{equation}
The exact shifted Dirichlet solution $x^+$ on $U\cup B$ is
\begin{align}
 x_B^+&=\mathcal S_B^{-1}q_B,
 \label{eq:lazy-schur-new-block}\\
 x_U^+&=x_U-A^{-1}Ex_B^+.
 \label{eq:lazy-schur-old-block}
\end{align}
For every remaining inactive block
$C\subseteq V\setminus(U\cup B)$,
\begin{align}
 e_\ell(x^+)_C
 &=e_\ell(x)_C-\mathcal S_{CB}x_B^+,
 \label{eq:lazy-schur-demand-update}\\
 \mathcal S_{CB}
 &:=Q_{CB}-Q_{CU}A^{-1}E\leq0.
 \label{eq:lazy-schur-cross-block}
\end{align}
Moreover, $\mathcal S_B$ is Stieltjes,
$\mathcal S_B^{-1}\geq0$, and $-A^{-1}E\geq0$.  In particular, if
$q_B>0$, then the new block is positive, the old coordinates only increase,
and every remaining violation demand can only increase.
\end{lemma}

\begin{proof}
The enlarged active equations are
\begin{equation*}
 \begin{bmatrix}A&E\\E^T&Q_{BB}\end{bmatrix}
 \begin{bmatrix}x_U^+\\x_B^+\end{bmatrix}
 =
 \begin{bmatrix}\ell_U\\\ell_B\end{bmatrix}.
\end{equation*}
The first row and $Ax_U=\ell_U$ give
\cref{eq:lazy-schur-old-block}.  Substitution in the second row gives
\cref{eq:lazy-schur-new-block}.  For an inactive block $C$,
\begin{align*}
 e_\ell(x^+)_C
 &=\ell_C-Q_{CU}x_U^+-Q_{CB}x_B^+\\
 &=e_\ell(x)_C-
   \bigl(Q_{CB}-Q_{CU}A^{-1}E\bigr)x_B^+,
\end{align*}
which is \cref{eq:lazy-schur-demand-update}.  Principal Schur complements of
a Stieltjes matrix are Stieltjes.  Their off-diagonal blocks are
nonpositive, and inverse positivity gives all remaining signs.  Since the
diagonal of $\mathcal S_B^{-1}$ is positive, $q_B>0$ also gives
$x_B^+>0$ coordinatewise.
\end{proof}

The dense response $-A^{-1}E$ need not be written into every old coordinate.
It may be stored as an elimination response, while
\cref{eq:lazy-schur-demand-update} updates only the boundary event
representation.  Old coordinates are recovered by one reverse substitution
after the final event.  This changes the implementation question from
\emph{how many old regions were settled?} to \emph{how compactly can the
signed Schur block-columns be represented and queried?}

\begin{theorem}[The fan cascade has an output-linear exact Schur trace]
\label{thm:fan-lazy-trace}
Use the parameters of \cref{thm:fan-cascade}, and consider its first $L$
singleton batches.  There is an exact implementation that
\begin{enumerate}[label=(\roman*),leftmargin=*]
 \item produces the identical batches
 $B_0=\{v_1\},\ldots,B_{L-1}=\{v_L\}$;
 \item represents every intermediate shifted Dirichlet solution exactly; and
 \item recovers $x^L$ on $U_L=\{o,v_1,\ldots,v_L\}$ exactly
\end{enumerate}
using $O(m+L)=O(\operatorname{cvol}(U_L))$ degree work and arithmetic.
In contrast, materializing and rescanning every fresh settlement through
$U_L$ costs $\Omega(mL)$.
\end{theorem}

\begin{proof}
Scan the degree-$m$ root list once.  Its degree metadata identifies the unique
degree-two neighbor $v_1$; thereafter one scan of each admitted $v_j$
identifies the next path neighbor $v_{j+1}$.  At stage $j\geq1$, retain the
separator
\begin{equation}
 Z_j:=\{o,v_j\},
 \qquad
 I_j:=U_j\setminus Z_j,
 \label{eq:fan-separator}
\end{equation}
and the exact reduced system obtained after eliminating $I_j$:
\begin{align}
 \widehat Q_j
 &:=
 Q_{Z_jZ_j}
 -Q_{Z_jI_j}Q_{I_jI_j}^{-1}Q_{I_jZ_j},
 \label{eq:fan-reduced-matrix}\\
 \widehat\ell_j
 &:=
 (\ell_\rho)_{Z_j}
 -Q_{Z_jI_j}Q_{I_jI_j}^{-1}(\ell_\rho)_{I_j}.
 \label{eq:fan-reduced-load}
\end{align}
Then
\begin{equation}
 x^j_{Z_j}=\widehat Q_j^{-1}\widehat\ell_j.
 \label{eq:fan-separator-solution}
\end{equation}
For $j=0$, the retained system is the scalar root equation.

The next vertex $v_{j+1}$ is adjacent to the old active region only through
$Z_j$.  Append it to the reduced system, producing a matrix of order at most
three, and eliminate $v_j$ to obtain
$(\widehat Q_{j+1},\widehat\ell_{j+1})$.  Store the eliminated pivot and its
two multipliers for the final reverse pass.  Each append, elimination, and
separator solve uses $O(1)$ arithmetic, and scanning the newly exposed fan
vertex costs at most three.

It remains to show that no repeated root scan is needed to reproduce the
gate.  For a degree-$d$ root neighbor with no active path neighbor, define
\begin{equation}
 \eta_d^{(j)}
 :=
 -\alpha\rho\sqrt d
 +\frac{1-\alpha}{2\sqrt{dm}}\,x_o^j.
 \label{eq:fan-live-remote-demand}
\end{equation}
Every remote $v_i$ with $i\geq j+2$ has the identical exact violation demand
$\eta_3^{(j)}$.  The unique next path vertex has exact demand
\begin{equation}
 e_{\ell_\rho}(x^j)_{v_{j+1}}
 =
 \begin{cases}
  \eta_2^{(0)},&j=0,\\
  \eta_3^{(1)}+h_1x_{v_1}^1,&j=1,\\
  \eta_3^{(j)}+hx_{v_j}^j,&j\geq2,
 \end{cases}
 \label{eq:fan-live-next-demand}
\end{equation}
where $h_1$ and $h$ are from
\cref{eq:fan-first-coupling,eq:fan-core-coupling}.  Hence the exact gate
decision at each of the first $L$ stages uses only the two separator
coordinates and two scalar comparisons.  The strict inequalities proved in
\cref{thm:fan-cascade} therefore reproduce the same singleton batches.

The initial root scan costs $m$.  The $L$ local scans, constant-size Schur
updates, and final reverse recovery cost $O(L)$.  Since $m>L$ and
\begin{equation*}
 \operatorname{cvol}(U_L)=m+4L,
\end{equation*}
this is output-linear.  Fresh settlement, by contrast, charges at least the
root's $m+1$ work units at each of the $L+1$ stages, giving
$\Omega(mL)$.
\end{proof}
